% Generated draft: requires mathematical and source review.
\documentclass[11pt]{article}
\usepackage[margin=25mm]{geometry}
\usepackage{fontspec}
\setmainfont{Latin Modern Roman}
\usepackage{amsmath,amssymb,mathtools,unicode-math}
\setmathfont{Latin Modern Math}
\usepackage{enumitem,needspace,xurl,hyperref}
\setlength{\parindent}{0pt}
\setlength{\parskip}{5pt}
\setlength{\emergencystretch}{3em}
\newcommand{\N}{\mathbb{N}}
\newcommand{\Z}{\mathbb{Z}}
\newcommand{\Q}{\mathbb{Q}}
\newcommand{\R}{\mathbb{R}}
\newcommand{\C}{\mathbb{C}}
\newcommand{\F}{\mathbb{F}}
\newcommand{\A}{\mathbb{A}}
\newcommand{\PP}{\mathbb{P}}
\newcommand{\E}{\mathbb{E}}
\newcommand{\eps}{\varepsilon}
\newcommand{\dd}{\,\mathrm{d}}
\begin{document}
\section*{UnsolvedMath ID: 800027}
\textbf{Source problem code:} catalogue entry
\subsection{Definitions and mathematical statement}
A semigroup is a nonempty set with an associative binary operation. A zero element of a semigroup \(S\) is an element \(0\) satisfying \(0s=s0=0\) for every \(s\in S\). A semigroup with zero is \(0\)-simple if \(S^2\ne\{0\}\) and its only nonempty two-sided ideals are \(\{0\}\) and \(S\). Here \(S^2=\{xy:x,y\in S\}\), and a nonempty subset \(I\subseteq S\) is a two-sided ideal if \(SI\cup IS\subseteq I\).

For a class \(\mathcal A\) of semigroups, the usual Birkhoff operators \(\mathbb H\), \(\mathbb S\), and \(\mathbb P\) take homomorphic images, subsemigroups, and direct products, respectively, with classes understood up to isomorphism. Products of operators are conventionally composed from right to left; for example,
\[
\mathbb H\mathbb P\mathbb S(\mathcal A)
=\mathbb H\bigl(\mathbb P(\mathbb S(\mathcal A))\bigr).
\]
The supplied excerpt does not specify whether \(\mathbb P\) permits arbitrary direct products or only finite direct products, nor whether homomorphisms and subsemigroups must preserve a distinguished zero. These conventions require confirmation.

\textbf{Problem (S. I. Kublanovsky; supplied Problem 28).} Let \(C_0\) be a class of finite \(0\)-simple semigroups. Is \(\mathbb K(C_0)\) decidable for each of the operators
\[
\mathbb K\in\bigl\{\mathbb H\mathbb P,\ \mathbb H\mathbb P\mathbb S,\ \mathbb S\mathbb H\mathbb P\mathbb S\bigr\}?
\]
The meaning of decidability is not specified in the supplied excerpt. If the intended inputs are finite semigroups represented by multiplication tables, the associated membership question is whether an algorithm, for the given \(C_0\) and chosen \(\mathbb K\), halts on every such input \(T\) and decides whether \(T\in\mathbb K(C_0)\). This is a conditional interpretation, not a supplied input convention. No effective description of \(C_0\), access model for \(C_0\), or requirement of uniformity in \(C_0\) is given. Consequently, the excerpt does not determine a unique fully specified decision problem.

The source additionally states that the answer for all other Birkhoff operators has been obtained in reference [Kubl6]. The supplied record gives neither that reference's bibliographic entry nor the precise results and hypotheses covered by this assertion.
\subsection{Short English statement}
For a class \(C_0\) of finite \(0\)-simple semigroups, ask whether membership in each of \(\mathbb H\mathbb P(C_0)\), \(\mathbb H\mathbb P\mathbb S(C_0)\), and \(\mathbb S\mathbb H\mathbb P\mathbb S(C_0)\) is decidable. The input domain, effective specification of \(C_0\), and operator conventions require clarification.
\subsection{Sources}
\begin{itemize}
\item Local UnsolvedMath registry: permanent ID \(800027\), source code \texttt{catalogue entry}, title \emph{The Lyapin's Notebook --- Problem 28}; supplied upstream revision \texttt{d371ee917ac9647fbb06cd6fc529deec6f3dda82}. The registry labels the problem open; this is supplied status metadata, not an independently verified assessment.
\item Supplied source list: Bershadsky, Kublanovsky, Mashevitzky, \emph{The Lyapin's Notebook: Unsolved Problems in Semigroup Theory} (2026), Problem 28, \url{https://arxiv.org/abs/2604.04763}. The problem text attributes the question to S. I. Kublanovsky; the registry's proposer field lists Bershadsky, Kublanovsky, Mashevitzky and records the year 2026.
\item Supplied provenance records access on 2026-10-16 and extraction by \texttt{arxiv\_latex}. These are supplied provenance details; no independent source visit or verification is claimed.
\item Reference [Kubl6] is cited in the supplied statement for the assertion about all other Birkhoff operators. Its bibliographic details and the applicable hypotheses are absent from the supplied record.
\end{itemize}
\end{document}
